%!TEX program = xelatex
%!TEX encoding = UTF-8
%Definições de Ano, Semestre, número e data da avaliação
\def\numeroteste{1}
\def\modulo{3}
\def\dataavaliacao{16/08/2024}

\documentclass[12pt]{exam}

\usepackage{caption}
\usepackage{amssymb}
\usepackage{amsmath,amsfonts,amsthm,amstext}
\usepackage[brazil]{babel}
% \usepackage[latin1]{inputenc}
\usepackage{graphicx}
\graphicspath{{/home/jfreitas/GitHub_Repos/Algebra-1/Pictures/}{D:/OneDrive - unb.br/imagens-latex/}}
\usepackage{enumitem}
\usepackage{multicol}
\usepackage{answers}
\usepackage{tikz,ifthen}
\usetikzlibrary{lindenmayersystems}
\usetikzlibrary[shadings]
\def\ano{2024}
\def\semestre{1}
\def\disciplina{Álgebra 1}
\def\nomeabreviado{Álgebra 1}
\def\turma{1}

\newcommand{\im}{{\rm Im\,}}
\newcommand{\dlim}[2]{\displaystyle\lim_{#1\rightarrow #2}}
\newcommand{\minf}{+\infty}
\newcommand{\ninf}{-\infty}
\newcommand{\cp}[1]{\mathbb{#1}}
\newcommand{\sub}{\subseteq}
\newcommand{\n}{\mathbb{N}}
\newcommand{\z}{\mathbb{Z}}
\newcommand{\rac}{\mathbb{Q}}
\newcommand{\real}{\mathbb{R}}
\newcommand{\complex}{\mathbb{C}}

\newcommand{\vesp}[1]{\vspace{ #1  cm}}

\newcommand{\compcent}[1]{\vcenter{\hbox{$#1\circ$}}}
\newcommand{\comp}{\mathbin{\mathchoice
        {\compcent\scriptstyle}{\compcent\scriptstyle}
        {\compcent\scriptscriptstyle}{\compcent\scriptscriptstyle}}}
\renewcommand{\sin}{{\rm sen\,}}
\renewcommand{\tan}{{\rm tg\,}}
\renewcommand{\csc}{{\rm cossec\,}}
\renewcommand{\cot}{{\rm cotg\,}}
\renewcommand{\sinh}{{\rm senh\,}}
\newcounter{exercicios}
\setcounter{exercicios}{0}
\newcommand{\questao}{
    \addtocounter{exercicios}{1}
    \noindent{\bf Questão \arabic{exercicios}: }}

\newcommand{\solucao}{
    \noindent{\bf Solução: }}

\newcommand{\resp}[1]{
    \noindent{\bf Exerc{\'\i}cio #1: }}

\extrafootheight[.25in]{.5in}
\footrule
\lfoot{Teste \numeroteste\ - Módulo \modulo\ - \nomeabreviado\ - Turma \turma\ - \semestre$^o$/\ano}
\cfoot{}
\rfoot{P\'agina \thepage\ de \numpages}

\begin{document}

    \begin{figure}[h]
        \begin{minipage}[c]{1.7cm}
            \includegraphics[width=1.7cm]{unb.pdf}
        \end{minipage}
        \hspace{0pt}
        \begin{minipage}[c]{4in}
            {Universidade de Brasília} \\
            {Departamento de Matemática}
        \end{minipage}
    \end{figure}
    \vesp{-1}
    \hrule
    \begin{center}
        {\Large\bf \disciplina\ - Turma \turma}  \\
        \vesp{0.2} {\large\bf Teste \numeroteste\ - Módulo\ \modulo\ -\ \dataavaliacao}
    \end{center}

    \noindent{\bf \underline{Nome: \hspace{9.8cm} Mat.:\qquad \
            \hspace{2.1cm}}}\\
    \vspace*{.01cm}

    \noindent{\bf \underline{Nome: \hspace{9.8cm} Mat.:\qquad \
            \hspace{2.1cm}}}

    \vspace{.4cm}


    \questao Seja $d \in \z$ e considere o subconjunto de $M_2(\z)$ dado por
    \[
    M_2^d(\z) = \left\{\begin{pmatrix} a & db \\ b & a \end{pmatrix} \mid a, b \in \z\right\}.
    \]
    Mostre que $M_2^d(\z)$ é um subanel de $M_2(\z)$.

    \solucao \noindent Primeiro observe que o elemento neutro do anel $M_2(\z)$ é
    \[
        0_{M_2(\z)} = \begin{bmatrix}
            0 & 0\\
            0 & 0
        \end{bmatrix}
    \]
    e é fácil ver que $0_{M_2(\z)} \in M_2^d(\z)$ e com isso, $M_2^d(\z) \ne \emptyset$.

    Agora sejam $x$, $y \in M_2^d(\z)$. Assim
    \begin{align*}
        x &= \begin{bmatrix}
            a & db\\b & a
        \end{bmatrix}\\
        y &= \begin{bmatrix}
            c & de\\e & c
        \end{bmatrix}
    \end{align*}
    e em $M_2(\z)$ o oposto de $y$ é:
    \[
        -y = \begin{bmatrix}
            -c & -de\\-e & -c
        \end{bmatrix}.
    \]

    Com isso temos
    \begin{align*}
        x + (-y) = \begin{bmatrix}
            a & db\\b & a
        \end{bmatrix} + \begin{bmatrix}
        -c & -de\\-e & -c
        \end{bmatrix} = \begin{bmatrix}
        a - c & d(b - e)\\b - e & a - c
        \end{bmatrix} \in M_2^d(\z).
    \end{align*}

    Além disso,
    \begin{align*}
        x \cdot y = \begin{bmatrix}
            a & db\\b & a
        \end{bmatrix} \cdot \begin{bmatrix}
            c & de\\e & c
        \end{bmatrix} = \begin{bmatrix}
            ac + dbe & ade + dbc\\bc + ae & bde + ac
        \end{bmatrix} = \begin{bmatrix}
            ac + dbe & d(ae + bc)\\ae + bc & ac + dbe
        \end{bmatrix} \in M_2^d(\z).
    \end{align*}

    Portanto $M_2^d(\z)$ é um subanel de $M_2(\z)$.

    \vspace{.4cm}

    \questao Seja $(A, +, \cdot)$ um anel comutativo. Se $I$ e $J$ são ideais de $A$, mostre que $I \cap J$ é um ideal de $A$.

    \solucao \noindent Como $I$ e $J$ são ideais de $A$, então $0_A \in I$ e $0_A \in J$, logo $0_A \in I \cap J$ e com isso, $I \cap J \ne \emptyset$.

    Agora sejam $x$, $y \in I \cap J$. Então $x$, $y \in I$ e $x$, $y \in J$. Mas $I$ e $J$ são ideais, logo $x + (-y) \in I$ e $x + (-y) \in J$. Logo $x + (-y) \in I \cap J$.

    Finalmente, seja $\alpha \in A$ e $x \in I \cap J$. Logo $x \in I$ e $x \in J$. Mas $I$ e $J$ são ideais, assim $\alpha x \in I$ e $\alpha y \in J$. Logo $\alpha x \in I \cap J$.

    Portanto $I \cap J$ é um ideal de $A$.
    \vspace{1cm}
\end{document}
